\documentclass[11pt]{article}
\usepackage[margin=1in]{geometry}
\usepackage{amsmath,amssymb,amsthm}
\usepackage[hidelinks]{hyperref}
\usepackage{xurl}
\newtheorem{lemma}{Lemma}
\newtheorem{theorem}[lemma]{Theorem}
\newtheorem{corollary}[lemma]{Corollary}
\newcommand{\rk}{\operatorname{rank}}
\title{An involutory rank-one obstruction without a preserved form}
\author{Written proof and exact algebra controls}
\date{17 November 2026}
\begin{document}\maketitle

\begin{theorem}\label{thm:373}
Over any field, the knot group of
$13a_{373}$ has no irreducible four-dimensional representation taking
meridians to nonidentity involutions of displacement rank one.
In characteristic different from two, it therefore has no irreducible
four-dimensional KCH representation with meridian eigenvalue $-1$.
No preserved symmetric or Hermitian form is assumed.
\end{theorem}

The literal PD, arc numbering and generating seed arcs $1,2,6,8$
are those of the retained seven-sheeted and universal-reflection
inputs. Their exact named-diagram identification is a dependency;
it is not inferred from this algebra calculation.
Write each meridian matrix as $R_i=I-u_if_i$, with $u_i,f_i$ nonzero.
Such a factorization exists by displacement rank one, and
\[
 R_i^2-I=(f_i(u_i)-2)u_if_i=0
\]
forces $f_i(u_i)=2$ over every field. In characteristic two this
value is zero and the matrix is a transvection; no nonzero norm or
diagonalizability assumption is used in the representation proof.
At a crossing $(o,a,b)$ the reflection pair propagates as
\[
 u_b=R_ou_a,\qquad f_b=f_aR_o.
\]
Pairs $(u,f)$ and $(tu,t^{-1}f)$ define the same matrix for every
nonzero scalar $t$. Two nonzero displacement factorizations define
the same matrix precisely when their pairs differ in this manner.

\begin{lemma}\label{lem:independent}
In an irreducible four-dimensional representation as in the theorem,
the four seed normals and the four seed covectors are both bases.
\end{lemma}
\begin{proof}
The span of the four seed normals is nonzero and invariant under
their four generating matrices, hence is the whole space. The
intersection of the four seed covector kernels is fixed pointwise
by the generators. Irreducibility makes this subspace zero: if it
were the whole space all generators would be identities. Thus the
four covectors span the dual space as well.
The normal-span argument bounds every irreducible representation
with rank-one meridian displacements by the number of generating
meridians, here four.
\end{proof}

\section{The first loop normalizes the directed pairings}

Suppose both families are independent. Take the four seed normals
as $e_0=u_1,e_1=u_2,e_2=u_6,e_3=u_8$, with corresponding covectors
$f_0,f_1,f_2,f_3$. Put $A_{ij}=f_i(e_j)$, so $A_{ii}=2$.
Let $a=A_{12}$ and $b=A_{21}$. The first three propagation steps
\[
                 (2,6,7),\ (7,2,3),\ (3,1,0)
\]
give
\[
\begin{aligned}
 u_3&=(1-ab)e_1+be_2,& f_3^{\rm arc}&=(1-ab)f_1+af_2,\\
 u_0&=e_0-su_3,& f_0^{\rm arc}&=f_0-tf_3^{\rm arc},
\end{aligned}
\]
where $s=(1-ab)A_{10}+aA_{20}$ and
$t=(1-ab)A_{01}+bA_{02}$.
Set $H=2-st$. The unused crossing $(0,1,2)$ implies
$R_0e_0=\kappa e_1$ and $f_0R_0=\kappa^{-1}f_1$
for some $\kappa\ne0$. Their coefficient vectors are
\[
 R_0e_0=(1-H)e_0+Hs(1-ab)e_1+Hsbe_2,
\quad
 f_0R_0=(1-H)f_0+Ht(1-ab)f_1+Htaf_2.
\]
The first equation gives $H=1$, hence $st=1$, then $b=0$
and $\kappa=s$. The second gives $a=0$; equivalently its
remaining term is $ta f_2=0$, while $t\ne0$ and $f_2\ne0$.
Thus $A_{12}=A_{21}=0$ and $A_{01}A_{10}=1$.
Replace $e_1$ by $se_1$ and $f_1$ by $s^{-1}f_1$.
Then both directed pairings $A_{01},A_{10}$ equal one.

The next steps $(1,3,4)$ and the preceding normalizations identify
the pairs on arcs $0,4,6,7$ with
\[
 (U,F)=(e_0-e_1,f_0-f_1),\quad(-U,-F),\quad
 (V,G)=(e_2,f_2),\quad(V,G),
\]
respectively. Put $(W,J)=(e_3,f_3)$.
Both triples $U,V,W$ and $F,G,J$ are independent; the latter
uses Lemma~\ref{lem:independent}. Each paired evaluation is two.

\section{A two-sided crossing contradiction}

Allow six unrestricted directed pairings
\[
 x=F(V),\ X=G(U),\ y=F(W),\ Y=J(U),\ z=G(W),\ Z=J(V).
\]
No transposition or conjugation relation is imposed on these values.
Set
\[
\begin{gathered}
 h=1-yY,\quad v=yX+zh,\quad C=Yx+Zh,\\
 w=(2+h)y-xv,\qquad Q=(2+h)Y-XC.
\end{gathered}
\]
The remaining propagation steps
$(8,4,5),(5,8,9),(6,9,10),(10,0,12),(4,10,11)$ give
\[
\begin{array}{c|c|c}
 i&u_i\text{ in }(U,V,W)&f_i\text{ in }(F,G,J)\\\hline
5&(-1,0,Y)&(-1,0,y)\\
9&(y,0,h)&(Y,0,h)\\
10&(y,-v,h)&(Y,-C,h)\\
11&(y-w,-v,h)&(Y-Q,-C,h)\\
12&(1-Qy,Qv,-Qh)&(1-wY,wC,-wh).
\end{array}
\]
These formulas, including every evaluation $f_i(u_i)=2$, are
checked as integer polynomial identities in the six directed variables.

The crossing $(11,7,8)$ says
$R_{11}V=\beta W$ and $GR_{11}=\beta^{-1}J$ for a nonzero
scalar $\beta$. Put $p=f_{11}(V)$ and $q=G(u_{11})$.
Independence of the root and covector triples gives
\[
\begin{array}{lll}
p(y-w)=0,&1+pv=0,&-ph=\beta,\\
q(Y-Q)=0,&1+qC=0,&-qh=\beta^{-1}.
\end{array}
\]
Thus $p,q,h$ are nonzero and
\[
             w=y,\quad Q=Y,\quad v=\beta^{-1}h,\quad C=\beta h.
\]
In particular the pair on arc12 simplifies to
\[
                 u_{12}=-hu_5+YvV,\qquad
                 f_{12}=-hf_5+yCG.
\]
The crossing $(12,5,6)$ says
$R_{12}u_5=\alpha V$ and $f_5R_{12}=\alpha^{-1}G$
for a nonzero scalar $\alpha$.
Put $r=f_{12}(u_5)$ and $s=f_5(u_{12})$.
The independent pairs $(u_5,V)$ and $(f_5,G)$ give
\[
           1+rh=0,\quad-rYv=\alpha,\qquad
           1+sh=0,\quad-syC=\alpha^{-1}.
\]
It follows that $\alpha=Y\beta^{-1}$ and
$\alpha^{-1}=y\beta$. Multiplication forces $yY=1$, so $h=0$,
contrary to $-ph=\beta\ne0$.
This proves Theorem~\ref{thm:373}.

\section{The other three reduced systems}

Consider the other three reduced eight-crossing systems for
$13a_{375},13a_{391},13a_{392}$ in the retained universal-reflection
input. The argument here uses their literal crossing relations, not
the preserved-form theorem. Both triples of seed roots and covectors
are again denoted $(U,V,W)$ and $(F,G,J)$, with six unrestricted
directed pairings as above.

For $13a_{375}$ and $13a_{391}$ let $(T,L)=R_V(W,J)$ and
$(A,H)=R_U(T,L)$, where this notation means simultaneous root and
covector conjugation. Put $t=F(T)$ and $\tau=L(U)$. Then
$A=T-tU$, $H=L-\tau F$.
For $13a_{375}$, the trace is
\[
 (6,9,10),(10,7,8),(4,10,11),(11,4,5),(8,11,12).
\]
It gives $(u_{10},f_{10})=(T,L)$, $(u_{11},f_{11})=(A,H)$,
$R_8=R_TR_VR_T$, and
\[
 u_5=(t\tau-1)U-\tau T,\qquad
 f_5=(t\tau-1)F-tL.
\]
The $U$-coefficient of $u_{12}$ is $-t$, and the $F$-coefficient
of $f_{12}$ is $-\tau$, since the pair on arc8 uses only $V,W,G,J$.
The unused crossing $(9,12,0)$ gives
$R_Wu_{12}=\gamma U$ and $f_{12}R_W=\gamma^{-1}F$.
Its two coefficients therefore force
$t=-\gamma$, $\tau=-\gamma^{-1}$, hence $t\tau=1$.
The pair on arc5 is now $(-\tau T,-tL)$, so $R_5=R_T$.
The unused crossing $(5,8,9)$ would imply
$\alpha W=R_5u_8=R_TR_TV=V$, a contradiction.

For $13a_{391}$ the trace is
\[
 (9,6,5),(6,9,10),(4,10,11),(5,11,12),(11,7,8).
\]
Here $(u_5,f_5)=R_W(V,G)$, $(u_{11},f_{11})=(A,H)$,
$(u_{12},f_{12})=R_5(A,H)$ and $(u_8,f_8)=R_A(V,G)$.
Again the coefficients on $U,F$ of the arc12 pair are $-t,-\tau$.
The unused crossing $(10,12,0)$ gives
\[
 u_{12}=\gamma(U-\tau T),\qquad
 f_{12}=\gamma^{-1}(F-tL).
\]
Thus $t=-\gamma$ and $\tau=-\gamma^{-1}$, and the displayed
pair equals $(A,H)$. Hence $R_{12}=R_A$.
The unused crossing $(12,8,9)$ now gives $V=\beta W$, impossible.

For $13a_{392}$ put $h=1-zZ$,
$t=zx+hy$, $\tau=ZX+hY$,
$k=(2+h)z-tX$ and $K=(2+h)Z-\tau x$.
The trace
\[
 (12,6,5),(5,12,11),(4,11,10),(6,10,9),(10,7,8)
\]
gives the paired coefficient vectors
\[
\begin{array}{c|c|c}
i&u_i\text{ in }(U,V,W)&f_i\text{ in }(F,G,J)\\\hline
11&(0,z,h)&(0,Z,h)\\
10&(-t,z,h)&(-\tau,Z,h)\\
9&(-t,z-k,h)&(-\tau,Z-K,h)\\
8&(Kt,1-Kz,-Kh)&(k\tau,1-kZ,-kh).
\end{array}
\]
The crossing $(9,12,0)$, on putting $p=f_9(W)$ and $q=J(u_9)$,
forces
\[
 pt=\gamma,\quad p(z-k)=0,\quad ph=1,\qquad
 q\tau=\gamma^{-1},\quad q(Z-K)=0,\quad qh=1.
\]
Thus $p,q,t,\tau,h$ are nonzero, $k=z$ and $K=Z$.
The crossing $(11,8,9)$ has root and covector phases $\beta,\beta^{-1}$.
Since the arc11 pair has zero coefficients on $U,F$, comparison of
those coefficients gives $Kt=-\beta t$ and
$k\tau=-\beta^{-1}\tau$. Hence $Z=-\beta$, $z=-\beta^{-1}$,
contradicting $h=1-zZ\ne0$.

In each of these four reduced systems, a three-seed trace propagates
all eight root directions and all eight covector directions. It follows
that they cannot simultaneously have root span three and covector
span three: the three chosen seed pairs would give the independent
triples used in the contradictions above. This is a two-sided
statement; it asserts no bound on just one span in a reducible case.

\section{Paired folds and the sixteen diagrams}

The local folds used to transfer these systems also work without a
form. Here are the needed arguments, with arbitrary nonzero closing
phase $\kappa$ in place of a sign. Take independent roots $A,E,C$
and independent covectors $F,G,J$. Write
\[
 a=F(E),\quad a'=G(A),\quad b=F(C),\quad b'=J(A),\quad
 c=G(C),\quad c'=J(E).
\]
In local typeI, the five crossings are
$(0,1,2),(j,2,3),(3,0,1),(1,3,4),(2,j,j')$,
where the pairs on arcs $0,1,j$ are $(A,F),(E,G),(C,J)$.
Set $t=c'-ab'$ and $T=c-a'b$. The arc3 pair is
\[
       (E-aA-tC,\ G-a'F-TJ).
\]
Put $H=-a'-Tb'$ and $K=-a-tb$.
The root and covector coefficients in the closing relation are
\[
 (1+aH,-H,tH)=(0,\kappa,0),\qquad
 (1+a'K,-K,TK)=(0,\kappa^{-1},0).
\]
Thus $t=T=0$, $aa'=1$. The arc4 pair is
$(-aA,-a'F)$ and the arc$j'$ pair equals $(C,J)$.
Consequently the full matrices satisfy $R_4=R_0$, $R_{j'}=R_j$.

In typeII replace the first two crossings by $(j,1,2),(0,2,3)$.
Set $t=a-c'b$, $T=a'-cb'$. The arc3 pair is
$(E-c'C-tA,G-cJ-TF)$.
The closing coefficient vectors are
\[
 (1-tT,T,-c'T)=(0,\kappa,0),\qquad
 (1-tT,t,-ct)=(0,\kappa^{-1},0).
\]
They force $c=c'=0$ and $aa'=tT=1$, giving the same full-matrix
identifications. These raw formulas and normalized paired folds are
checked with all six directed variables independent.

For the alternate old-four local loop
$(3,1,0),(1,3,4),(0,3,2),(2,6,7),(7,1,2)$,
write $A=u_1,E=u_3,C=u_6$. The pair on arc2 is
$(aA+(1-aa')E,a'F+(1-aa')G)$.
The pair on arc7 has root coefficient one on $C$.
The closing root equation therefore first forces its evaluation on
$A$ to vanish, and then $A=\kappa u_2$. Hence $aa'=1$.
Its evaluation condition gives $b'=0$. The closing covector equation
then forces the other evaluation to vanish, giving $b=0$.
Thus $R_7=R_6$ and $R_4=R_0$, with no symmetric-pairing hypothesis.

All seed changes in the retained sixteen-diagram transfer are
compositions of the following elementary operation. Conjugate one
seed pair by a reflection whose root and covector lie in the spans
of the other three retained pairs. The replaced root has coefficient
one on its previous root; the replaced covector likewise has
coefficient one on its previous covector. Both changes have
determinant one. This applies as well when one first changes the
other member of a pair and then performs a second such operation.
The exact word identities and changed seed lists are retained in
the input; they are diagrammatic conjugation identities and use no
form. For $13a_{391}$ use $u_6=R_9u_5$. For $13a_{392}$ first use
$u_{12}=R_5u_{11}$, then $u_6=R_{12}u_5$.
For $13a_{3264}$ and $13a_{3276}$ put $u_0=R_3u_1$, then $u_{12}=R_Qu_0$ and
$u_6=R_{12}u_5$, with $Q=10,9$ respectively.
For $13a_{3266}$ first use $u_9=R_5u_8$, then $u_{12}=R_9u_0$ and
$u_6=R_{12}u_5$. For $13a_{3275}$ use $u_6=R_9u_5$.
For the first two A6 rows below use $u_6=R_Qu_5$ with $Q=9,11$;
the third row requires no change, and the last uses $u_8=R_5u_7$.
The corresponding covectors are always propagated contragrediently.
The checker reconstructs these literal crossing words and computes
both change-of-basis determinants as the constant integer polynomial
one in twelve independent directed pairings. Thus independence of
both four-seed families is preserved in every listed change.

For the original remaining four knots, the folded eight-crossing
systems are identical label for label to those for373,375,391,392,
respectively, in the order3264,3266,3275,3276.
For the eight six-sheeted alternating cases, the exact bijections are
\[
\begin{array}{c|c|c|c}
\text{new indices}&\text{new labels}&\text{old labels}&\text{old index}\\\hline
304,3232&(0,5,6,8,9,10,11,12)&(9,10,11,12,0,5,6,8)&375\\
362,3259&(0,5,6,8,9,10,11,12)&(10,11,12,0,5,6,8,9)&375\\
623,3374&(0,5,6,7,8,10,11,12)&(8,6,5,0,12,11,10,9)&373\\
629,3375&(0,5,6,7,8,10,11,12)&(5,0,12,11,10,9,8,6)&373.
\end{array}
\]
Every index here is a $13a$ knot. The first index in each pair uses
typeI and the second typeII. The literal PD crossing sets verify
all twelve transfers directly, allowing only exchange of under-arcs.
The surviving reduced directions include three independent roots
and three independent covectors, since they are subsets or elementary
changes of the four-seed bases. This contradicts the two-sided
reduced-system conclusion above.

\begin{theorem}\label{thm:sixteen}
For every knot in
\[
\begin{aligned}
\mathcal K_{16}=\{&13a_{373},13a_{375},13a_{391},13a_{392},\\
 &13a_{3264},13a_{3266},13a_{3275},13a_{3276},\\
 &13a_{304},13a_{362},13a_{623},13a_{629},\\
 &13a_{3232},13a_{3259},13a_{3374},13a_{3375}\},
\end{aligned}
\]
every irreducible representation taking meridians to nonidentity
involutions of displacement rank one has dimension at most three,
over every field. No preserved form is assumed. In characteristic
different from two this includes every irreducible KCH representation
with meridian eigenvalue $-1$; in characteristic two it is a
transvection obstruction, rather than a KCH eigenvalue assertion.
\end{theorem}

\section{The exact augmentation consequence and primary comparison}

\begin{corollary}\label{cor:augmentation}
Let $K\in\mathcal K_{16}$ and let $C_K$ be the $U=1$ cord algebra in
the conventions of Cornwell's author version arXiv:1310.7526v3.
Every augmentation $\epsilon:C_K\to k$ with $\epsilon(\mu)=-1$ has rank
at most three, over every field $k$ of characteristic different
from two. No condition on the nonzero longitude value
$\epsilon(\lambda)$ is imposed.
\end{corollary}
\begin{proof}
Cornwell's Theorem1.2 constructs an irreducible KCH representation
from every augmentation with $\epsilon(\mu)\ne1$; Definition1.8
defines its rank to be the representation dimension, and Lemma3.14
identifies this with the cord-matrix rank \cite{cor17}.
Proposition2.12 identifies $\epsilon(\mu)$ with the distinguished
meridian eigenvalue and $\epsilon(\lambda)$ with its longitude
eigenvalue. The hypothesis $\epsilon(\mu)=-1$ is therefore exactly
the involutory case above. Lemma~\ref{lem:independent} excludes
dimensions greater than four, and Theorem~\ref{thm:sixteen} excludes
dimension four. These correspondence statements are explicit imports;
their original proofs are not independently re-established here.
\end{proof}

The specialization $\mu=-1$ does not automatically make an augmentation
reflective. Cornwell's arXiv:1509.04962v1, printed p.22 in the proof
of Corollary6.11, requires both $\epsilon(\mu)=-1$ and
$\epsilon(c_{ij})=\epsilon(c_{ji})$ \cite{cor15}.
Its printed p.8, after Theorem3.6, notes that every nontrivial knot
has a rank-two complex augmentation with meridian value $-1$, whereas
Corollary5.3 on printed p.15 gives determinant-one knots without
rank-two complex reflective augmentations. Thus dropping the symmetry
hypothesis is a substantive change of scope. This comparison is not
an assertion of worldwide priority for the particular obstruction.

\section{Exact certificate and scope}

The standard-library program \path{verify_involutory_kch_obstruction.py}
checks the raw first-loop formulas, the five reduced paired traces
and all ten root/covector expressions. It also verifies a constant-one
identity in the integer Laurent ring with $\alpha,\beta$ inverted,
using eleven necessary closure equations and26 multiplier monomials.
The expanded certificate is retained as \texttt{unit\_identity} in
\path{involutory_kch_obstruction_inputs.json};
a changed multiplier is rejected. No finite-field sample, numerical
solution or untrusted Gr\"obner output supplies the contradiction.
The same standalone checker verifies the other
three paired systems, both raw six-variable A6 local folds, the
alternate old-four fold and all twelve literal PD incidence transfers.

The independence of both families is essential to this proof.
It does not assert that every reducible rank-one involutory
representation has fixed-space codimension at most three.
It gives no bound for noninvolutory meridian images and no upper
bound on augmentation rank over all meridian parameters.
The augmentation consequence uses the exact published correspondence;
that correspondence is not a premise of the representation theorem.
All checks and proofs here are internal and have not received
formal verification or external specialist review.
\begin{thebibliography}{9}
\bibitem{cor17} C.R.Cornwell, KCH representations, augmentations,
and $A$-polynomials, J.Symplectic Geom.15 (2017),983--1017.
Retained author version arXiv:1310.7526v3, Theorem1.2 (p.1;
proof pp.12--13), Definition1.8 (p.3), Proposition2.12 (p.9),
Lemma3.14 (p.13).
\url{https://arxiv.org/abs/1310.7526v3}.
\bibitem{cor15} C.R.Cornwell, Character varieties of knot complements
and branched double-covers via the cord ring, author preprint
arXiv:1509.04962v1 (2015), Theorem3.6 and following paragraph (p.8),
Corollary5.3 (p.15), and proof of Corollary6.11 (p.22).
\url{https://arxiv.org/abs/1509.04962v1}.
\end{thebibliography}
\end{document}
